\chapter{Datasets and Preprocessing}
\label{ch:data}

Three public datasets are used: YooChoose for session-based recommendation, and
Instacart and Dunnhumby for next-basket recommendation. All three are downloaded
from Kaggle by the notebook. All preprocessing follows published protocols so
that the results can be compared with prior work, and every statistic in this
chapter is computed and exported by the notebook.

\section{YooChoose (Session-Based)}

YooChoose is the click log of the ACM RecSys Challenge 2015 \cite{benshimon2015},
collected from a European e-commerce site over six months in 2014. Each row
contains a session ID, a timestamp, an item ID and a category; the click file has
33,003,944 rows.

\subsection{Protocol}

The preprocessing follows GRU4Rec \cite{hidasi2016} and the yoochoose1/64 split
used by NARM, STAMP and SR-GNN \cite{li2017,liu2018,wu2019}:

\begin{enumerate}
    \item Remove sessions with a single click, then items that occur in fewer than
          5 sessions, then sessions that have become shorter than 2 clicks.
    \item \textbf{Test set}: all sessions that end during the last day of the log.
    \item \textbf{Training data}: the most recent 1/64 of the remaining sessions
          (ordered by start time). This is the common yoochoose1/64 setting; the
          UPES profile also runs yoochoose1/4.
    \item \textbf{Validation set}: the sessions of the last day of that training
          data are held out for early stopping and model selection.
    \item Validation and test keep only items that occur in the final training
          data; sessions shorter than 2 clicks after this filter are removed.
\end{enumerate}

Every click after the first in a validation or test session is a prediction
target, so a session of length $n$ contributes $n-1$ predictions. Because the
split is by session and by time, no session contributes to more than one split.
The notebook asserts this, and that every training session ends before the
validation window and every validation session before the test window. This
removes the leakage of the first draft, where prefixes of the same session were
shuffled into both training and test sets.

\generated{dataset_yoochoose.tex}

The resulting split closely matches the one reported by SR-GNN for
yoochoose1/64 (55,898 test predictions and 16,766 items in \cite{wu2019}).

\begin{figure}[H]
\centering
\genfigure{eda.pdf}{width=\textwidth}
\caption{Left: clicks per session in the raw YooChoose log (log scale). Centre:
item popularity by rank (log-log), showing the long tail. Right: number of
baskets per user in the two next-basket datasets.}
\label{fig:eda}
\end{figure}

\section{Instacart (Next-Basket)}

The Instacart Market Basket Analysis data contain over three million grocery
orders from 206,209 users, with 49,688 products in 134 aisles. Each user has
between 4 and 100 orders. Kaggle's \texttt{eval\_set} column marks every user's
final order as either \emph{train} (products given) or \emph{test} (products
hidden for the competition).

\subsection{Protocol}

Following van Maasakkers et al.\ \cite{vanmaasakkers2023}:

\begin{enumerate}
    \item Drop the hidden Kaggle \emph{test} orders; each user's last remaining
          order is the target basket and all earlier orders are the history.
    \item \textbf{Product clustering.} Rare products (fewer than 500 purchases in
          the training history) are merged within their aisle. The authors'
          R script merges greedily by cosine similarity of the products' purchase
          vectors; we re-implemented it in Python, merging the rarest cluster with
          its most similar rare cluster in the same aisle until every cluster
          reaches 500 purchases or the aisle has no rare products left. Frequent
          products keep their own index. Products that appear only in target
          baskets are removed.
    \item Users are split at random into equal validation and test groups. The
          target basket of a validation user is used for early stopping; the
          target basket of a test user for the final evaluation. The histories of
          all users are used for training.
\end{enumerate}

Our clustering yields 12,058 product groups, against 9,407 in the original
paper, because our implementation only merges rare products with each other,
whereas the original evidently also merged more aggressively. On the Mac profile a random 10\% of users is used, as in the
NBR re-evaluation study \cite{li2023}; the GPU profiles use all users.

\section{Dunnhumby: The Complete Journey (Next-Basket)}

The Complete Journey contains two years of transactions of 2,500 households at a
retailer: 2,595,732 transaction rows over 92,339 products, with the household,
basket, day and time of each purchase.

\subsection{Protocol}

Following van Maasakkers et al.\ \cite{vanmaasakkers2023}:

\begin{enumerate}
    \item Keep products bought at least 50 times and households with at least 3
          baskets.
    \item Order each household's baskets by day and transaction time and keep the
          \emph{100 most recent} baskets. (The published script keeps the first
          100 baskets, which discards the most recent history of 37\% of the
          households; we keep the most recent ones.)
    \item The last basket of each household is the target; households are split
          at random into equal validation and test groups.
\end{enumerate}

After preprocessing the data contain 2,478 households, 10,370 products, 163,937
baskets and 66.16 baskets per household. These figures are identical to those
reported by van Maasakkers et al.

Note that several next-basket benchmarks, including TIFU-KNN \cite{hu2020} and
the NBR re-evaluation \cite{li2023}, use a different Dunnhumby dataset (``Let's
Get Sort-of-Real''), so their Dunnhumby results are not comparable with ours.

\generated{dataset_baskets.tex}

\section{Repeat Behaviour}

The ``Repeat ratio'' row of Table~\ref{tab:basket_stats} gives the average share
of a target basket that the customer had bought before. It is about 0.60 on
Instacart and 0.53 on Dunnhumby, in line with the 60\% and 51\% reported by van
Maasakkers et al. Repeat items are much easier to predict than new ones
\cite{li2023}, which is why the results in Chapter~\ref{ch:results} are also
reported separately for repeat and explore items.
